\documentclass[14pt]{extarticle}

\usepackage{labstyle}
\usepackage{physics}

\begin{document}
	\textbf{Цель работы.} Изучение основных процессов происходящих при прохождении сигналов через радиотехнические цепи с нелинейными элементом, экспериментальное исследование полупроводникового преобразователя частоты и амплитудного диодного детектора.
	
	\textbf{Назначение устройства.} Исследуемые в` лабораторной работе устройства используется для нелинейного преобразования сигналов. Преобразователь частоты -- для перемещения спектра сигнала по шкале частот без изменения характеристика сигнала, т.е. соотношений между компонентами спектра. Например, для модулированных колебаний это означает изменение (повышение или понижение) несущей частоты с сохранением вида модуляции и закона изменения модулированного параметра. Амплитудный детектор -- для преобразования модулированного колебания в низкочастотное колебание, соответствующее модулирующему сигналу.
	
	Основные радиотехнические преобразования сигналов осуществляются с помощью нелинейных цепей. Поэтому свойства нелинейных элементов и цепей являются фундаментом для теории большинства радиотехнических устройств. 
	
	Нелинейными называются цепи, описываемые нелинейными дифференциальными уравнениями вида \cref{eq:th:1}
	\begin{align} \label{eq:th:1}
		a_0 \dv[n]{y}{t} + a_1 \dv[n-1]{y}{t} + \dots = b_0 \dv[m]{x}{t} + b_1 \dv[m-1]{x}{t} + b_m x,
	\end{align}  
	в которых хотя бы один из коэффициентов $a_i$ является функцией $y$ или ее производных, либо один из коэффициентов $b_j$ -- функцией $x$ или ее производных:
	\[
	a_i = a_i \left(y, \dv{y}{t}, \dots\right); \qquad b_j = b_j\left(x, \dv{x}{t}, \dots\right).
	\] 
	
	Уравнение электрической цепи оказывается нелинейным в том случае, когда в схеме используются какие-либо нелинейные элементы, т.е. элементы параметры которых зависят от тока или напряжения. Нелинейными являются все электронные и полупроводниковые приборы. Особенности этих приборов определяются вольт-амперной характеристиками, т.е. зависимостями токов от приложенных напряжений.
	
	\subsection{Аппроксимация вольт-амперных характеристик}
	
	Математическое описание работы радиотехнической цепи начинается с составления уравнений, связывающих токи и напряжения в различных частях ее, в том числе и в нелинейных элементах. Для нелинейных элементов обычно известна графическая зависимость тока от напряжения (из справочника или эксперимента). 
\end{document}